\section{Introduction}

Let $R$ be a commutative integral domain of characteristic 0, and let $R\la X,Y\ra$ be
the free associative algebra generated over $R$ by two (non-commutative) 
letters $X$ and~$Y$.
For $u,v\in R\la X,Y\ra$, we shall write $[u,v]$
to denote the Lie bracket $uv-vu$.
In \cite{M37}, W.~Magnus introduced the associative subalgebra 
$S_X\subset R\la X,Y\ra$ generated by 
(what are called) the elements arising by elimination of $X$:
\begin{equation}
Y^{(0)}:=Y, \quad
Y^{(k+1)}:=[X,Y^{(k)}] \quad
(k=0,1,2,\dots),
\end{equation}
and showed that $S_X$ is freely generated by the $Y^{(k)}$
$(k=0,1,2,\dots)$.
Moreover, he derived that every element $Z$ of  $R\la X,Y\ra$ can be
written uniquely in the form
\begin{equation} \label{1st-elimination}
Z=\alpha_0 X^m +s_1 X^{m-1}+\cdots+s_m,
\end{equation}
where $\alpha_0\in R$, $s_1,\dots, s_m\in S_X$
(see \cite[Hilfssatz 2]{M37}, \cite[Lemma 5.6]{MKS}).
This observation is the first step in the construction 
of the basic Lie elements (an ordered basis of the free Lie algebra), 
which are obtained via repeated elimination, and
whose powered products in decreasing orders give 
the Poincar\'e-Birkoff-Witt
basis of the enveloping algebra~$R\la X,Y\ra$ (\cite[Theorem~5.8]{MKS}).
Apparently, this theory was historically a starting point 
toward subsequent developments of
finer constructions of free Lie algebra bases 
due to Lazard, Hall, Lyndon, Viennot and others
(see, e.g., \cite[Notes 4.5, 5.7]{Reu}). 

In this note, we however stay on the first step of elimination 
\eqref{1st-elimination} and 
look at combinatorial properties of a certain basis 
$\{\MM^{(\bk)}\}_{\bk\in \N_0^{(\infty)}}$ of $R\la X,Y\ra$
(to be called the Magnus polynomials below)
designed as follows: 

\begin{Notation} 
Let $\N_0$ denote the set of non-negative integers, and 
let 
$$
\N_0^{(\infty)}:=\bigcup_{d=0}^\infty
\left(
\prod_{k=1}^{d}\N_0\right)
\times \N_0 
$$ 
be the collection of finite sequences 
$
\bk=(k_1,\dots,k_d; k_\infty)
$
of non-negative integers equipped 
with a special last entry $k_\infty\in \N_0$. 
Here, we consider $(;k_\infty)$ also as elements of 
$\N_0^{(\infty)}$ coming from $d=0$.
For $\bk\in \N_0^{(\infty)}$, define
$|\bk|:=\sum_{i=1}^\infty k_i=k_1+\cdots+k_d+k_\infty$ 
(resp. $\dep(\bk):=d$),
and call it the size (resp. depth) of $\bk$.
\end{Notation}

\begin{Definition}[Magnus polynomial]
\label{defMagnus}
For $\bk=(k_1,\dots,k_d; k_\infty)\in\N_0^{(\infty)}$, define
$$
\MM^{(\bk)}:=Y^{(k_1)}\cdots Y^{(k_d)} \cdot X^{k_\infty} \in R\la X,Y\ra.
$$
We also set $\MM^{(;0)}=1$, and $\MM^{(;k)}=X^k$ for $k=1,2,\dots$.
Note that $\MM^{(k;0)}=Y^{(k)}$ for~$k\ge 0$.
\end{Definition}

\begin{Example} 
$\MM^{(1,0;2)}=Y^{(1)}Y^{(0)}X^2=(XY-YX)YX^2=XY^2X^2-YXYX^2$.
\end{Example}

It is not difficult to see that the Magnus polynomial $\MM^{(\bk)}\in R\la X,Y\ra$ 
is homogeneous
of bidegree $(|\bk|,\dep(\bk))$ in $X$ and $Y$.

The above mentioned Magnus expression \eqref{1st-elimination}
can then be rephrased as
\begin{equation} \label{Magnus}
Z=\sum_{\bk\in \N_0^{(\infty)}}
\alpha_{\bk}\, \MM^{(\bk)}
\end{equation}
with uniquely determined coefficients $\alpha_{\bk}\in R$ for any given 
$Z\in R\la X,Y\ra$.
In other words, the collection $\{\MM^{(\bk)}\mid \bk\in  \N_0^{(\infty)}\}$ 
forms an $R$-linear basis of $R\la X,Y\ra$.

Below in \S 2, we will construct another $R$-linear basis 
$\{\SS^{(\bk)}\mid \bk\in  \N_0^{(\infty)}\}$
(formed by what we call the ``demi-shuffle'' polynomials)
and show that 
$\{\MM^{(\bk)}\}_\bk$
and $\{\SS^{(\bk)}\}_\bk$
are dual to each other under the standard pairing
with respect to the monomials of~$R\la X,Y\ra$
(Theorem \ref{orthogonality}).
We then in \S 3 shortly generalize the duality to the case of
free associative algebras of more variables (Theorem 
\ref{orthogonality2}). 
In \S 4, we apply the formation of dual basis 
to derive a formula of Le--Murakami, Furusho type 
that expresses arbitrary coefficients of a group-like series 
$J\in R\lala X,Y\rara$ in terms of 
the ``regular'' coefficients of $J$
(Theorem \ref{LeMuFu}).  

\section{Demi-shuffle duals and array binomial coefficients}

Let $W$ be the subset of $R\la X,Y\ra$ formed by the monomials in $X,Y$ 
together with $1$, and call any element of $W$ a word. 
It is clear that $W$ forms a free monoid by the concatenation product
that restricts the multiplication of $R\la X,Y\ra$.
Each element of $R\la X,Y\ra$ is an $R$-linear combination of words in $W$.
For two elements $u,v\in R\la X,Y\ra$, define the standard pairing $\lp u,v\rp\in R$
so as to extend $R$-linearly the Kronecker symbol $\lp w,w'\rp:=\delta_{w}^{w'}\in\{0,1\}$
for words $w,w'\in W$.

\begin{Notation} We use the notation $w_\bk:=X^{k_1}Y\cdots X^{k_d}YX^{k_\infty}$
and call it the word associated to $\bk=(k_1,\dots,k_d; k_\infty)\in\N_0^{(\infty)}$.
The mapping $\bk\mapsto w_\bk$ gives a bijection from $\N_0^{(\infty)}$ onto $W$.
(Note that $w_{(;0)}=1$.)
The standard pairing $\lp w_\bk,w_{\bk'}\rp$ is equal to~$0$ or $1$ 
according to whether $\bk\ne \bk'$ or $\bk=\bk'$.
\end{Notation}


The purpose of this section is to describe the dual of the 
Magnus basis
$\{\MM^{(\bk)}\}_{\bk\in \N_0^{(\infty)}}$ with respect to the standard
pairing.

 
\begin{Definition}[Demi-shuffle polynomial]
\label{demi-shuffle}
For $\bk=(k_1,\dots,k_d; k_\infty)\in\N_0^{(\infty)}$, define
$$
\SS^{(\bk)}:=
(\cdots((X^{k_1}Y)\sha X^{k_2})Y)\sha\cdots)\sha  X^{k_d})Y)\sha
X^{k_\infty} \in R\la X,Y\ra,
$$
where $\sha$ denotes the usual shuffle product. 
We also set $\SS^{(;0)}=1$, and $\SS^{(;k)}=X^k$ for $k=1,2,\dots$.
Note that $\SS^{(k;0)}=X^kY$ for $k\ge 0$.
\end{Definition}

The construction of $\SS^{(\bk)}$ can be interpreted 
as forming the linear sum
of all words obtained from 
the word
$w_\bk=X^{k_1}Y\cdots X^{k_d}YX^{k_\infty}$
by consecutively applying ``left shuffles'' of the $X$ letters
and ``concatenations'' of the $Y$ letters in $w_\bk$.

\begin{Example}
Here are a few examples:
$\SS^{(0,1;0)}=(Y\sha X)Y=YXY+XYY$;
$\SS^{(1,1;0)}=((XY)\sha X)Y=XYXY+2XXYY$;
$\SS^{(1,0,1;0)}=(((XY)Y)\sha X)Y=XYYXY+XYXY^2+2X^2Y^3$.
Using the first identity, one can also compute
\begin{align*}
&\SS^{(0,1;1)}=((Y\sha X)Y)\sha X=(YXY+XYY)\sha X \\
&=(YXYX+2YXXY+XYXY)+(XYYX+XYXY+2XXYY) \\
&=2XXYY+2XYXY+XYYX+2YXXY+YXYX.
\end{align*}
\end{Example}

\begin{Theorem}[Duality] \label{orthogonality}
For $\bt, \bk\in \N_0^{(\infty)}$, we have
$$
\lp \SS^{(\bt)},\MM^{(\bk)}\rp=\delta_\bt^{\bk}.
$$
Here $\delta_\bt^{\bk}$ is the Kronecker symbol, i.e., designating $0$ or $1$
according to whether $\bt\ne\bk$ or $\bt=\bk$ respectively.
\end{Theorem}

Before going to the proof of the above theorem, we introduce the following
notation.

\begin{Definition}[Array binomial coefficient]
For $\bt,\bk\in  \N_0^{(\infty)}$ with $\dep(\bt)=\dep(\bk)$, 
$|\bt|=|\bk|$, define
\begin{equation} \label{ArrayBino}
\binom{\bt}{\bk}:=
\binom{t_1}{k_1}
\binom{t_1+t_2-k_1}{k_2}
\cdots
\binom{t_1+\cdots t_d-k_1-\cdots-k_{d-1}}{k_d},
\end{equation}
where $\bt=(t_1,\dots,t_d,t_\infty)$,
$\bk=(k_1,\dots,k_d,k_\infty)$.
We understand $\binom{\bt}{\bk}=1$ if
$\bt=\bk=(;N)$ for some $N\in \N_0$.
We set $\binom{\bt}{\bk}:=0$ if either $\dep(\bt)\ne\dep(\bk)$ or 
$|\bt|\ne|\bk|$ holds.
\end{Definition}

\begin{Remark} 
\label{rem2.6}
The special case $\binom{N,0,\dots,0;0}{k_1,k_2,\dots,k_d;k_\infty}$ 
is the same as the usual multinomial coefficient 
$\binom{N}{k_1,k_2,\dots,k_d,k_\infty}$ in combinatorics.
Note also that $\binom{\bt}{\bk}\ne 0$ implies $t_\infty\le k_\infty$,
as the last factor of $\binom{\bt}{\bk}$ could survive only when 
$(t_1+\cdots t_d-k_1-\cdots-k_{d-1})-k_d=k_\infty-t_\infty\ge 0$.
\end{Remark}

It turns out that the array binomial coefficients give
the expansion of $\SS^{(\bt)}$ as a linear sum of the monomials in $W$.
Recall that, for $\bt=(t_1,\dots,t_d,t_\infty)\in  \N_0^{(\infty)}$, 
$w_\bt$ denotes the word $X^{t_1}YX^{t_2}Y\cdots X^{t_d}YX^{t_\infty}\in W$.



\begin{Lemma}[Monomial expansion]
\label{dS-mono}
$$
\SS^{(\bk)}=\sum_{\bt \in\N_0^{(\infty)}}
\binom{\bt}{\bk} w_{\bt}.
$$
\end{Lemma}
\begin{proof}
Without loss of generality, it suffices to show 
$\la w_{\bt}, \SS^{(\bk)} \ra=\binom{\bt}{\bk}$
in the case $(N:=)\,|\bt|=|\bk|$ and $(d:=)\,\dep(\bt)=\dep(\bk)$.
The assertion is trivial when $d=0$, as then
$\bk=\bt=(;N)$, $\SS^{(\bk)}=X^N=w_\bt$ and $\binom{\bt}{\bk}=1$.
For $d>0$, we argue by induction on $d$.
Suppose $d=1$, $\bk=(k_1;k_\infty)$ and $\bt=(t_1;t_\infty)$. Then
$$
\SS^{(\bk)}=
(X^{k_1}Y)\sha X^{k_\infty}
=\sum_{i=0}^{k_\infty}(X^{k_1}\sha X^i)YX^{k_\infty-i}
=\sum_{i=0}^{k_\infty}\binom{k_1+i}{k_1}
X^{k_1+i}YX^{k_\infty-i}.
$$
Since $N=k_1+k_\infty=t_1+t_\infty$, we have
$\la  w_{\bt}, \SS^{(\bk)} \ra=\binom{k_1+k_\infty-t_\infty}{k_1}
=\binom{t_1}{k_1}$.
Suppose $d>1$ with $\bk=(k_1,\dots,k_d;k_\infty)$ and
$\bt=(t_1,\dots,t_d;t_\infty)$.
Write $\bk'=(k_1,\dots,k_{d-1};0)\in \N_0^{(\infty)}
$. Then
\begin{align*}
\SS^{(\bk)}
&=(((\SS^{(\bk')}Y)\sha X^{k_d})Y)\sha X^{k_\infty} 
\\
&=\sum_{i=0}^{k_\infty}
\SS^{(\bk')}
Y(X^{k_d}\sha X^i)YX^{k_\infty-i} 
\qquad (\text{associativity of $\sha$})
\\
&=\sum_{i=0}^{k_\infty}
\sum_{\bt'}\binom{\bt'}{\bk'}
\binom{k_d+i}{k_d}
w_{\bt'}
X^{k_d+i}YX^{k_\infty-i},
\end{align*}
where $\bt'=(t_1',\dots,t_{d-1}';t_d')\in \N_0^{(\infty)}$ runs over
those tuples 
with $t_1'+\cdots+t_d'=|\bk'|$ %k_1+\cdots+k_{d-1}$
so that $\SS^{(\bk')}$ is expressed as
$\sum_{\bt'}\binom{\bt'}{\bk'}w_{\bt'}$
by the induction hypothesis on $\dep(\bk')=d-1$.
The coefficient of $w_\bt$ in $\SS^{(\bk)}$
can be found in the above summand where
$k_\infty-i=t_\infty$,
$t_d'+k_d+i=t_d$ and
$t_s'=t_s$ ($s=1,\dots,d-1$), hence
$$
\la w_\bt, \SS^{(\bk)} \ra=
\binom{t_1}{k_1}
\cdots
\binom{t_1+\cdots t_{d-1}-k_1-\cdots-k_{d-2}}{k_{d-1}}
\cdot
\binom{k_d+k_\infty-t_\infty}{k_d}.
$$
Since $N=|\bk|=|\bt|$, we have
$k_d+k_\infty-t_\infty
=t_1+\cdots t_{d}-k_1-\cdots-k_{d-1}$.
This establishes the formula  
$\la w_\bt, \SS^{(\bk)} \ra=\binom{\bt}{\bk}$.
\end{proof}

\begin{Remark}
It would be worth noting that 
Lemma \ref{dS-mono} 
can be derived from
counting~$\la w_{\bt}, \SS^{(\bk)} \ra$ 
as the number of certain shufflings of letters in
$w_\bk=X^{k_1}Y\cdots X^{k_d}YX^{k_\infty}$
to produce 
$w_\bt=X^{t_1}Y\cdots X^{t_d}YX^{t_\infty}$.
Assume 
$|\bt|=|\bk|$ and $\dep(\bt)=\dep(\bk)$,
and consider letters $Y$ as partitions between 
groups of letters $X$ in $w_\bk$ and in $w_\bt$. 
Then~$\la w_{\bt}, \SS^{(\bk)} \ra$ is the number of ways
of moving some letters $X$ in $w_\bk$
to the left (beyond any number of $Y$'s)
to form the word $w_\bt$
without changing orders between~$X$'s from the same group
in $w_\bk$.
We count this number 
by enumerating branches of possibilities
for choosing
places of~$X$'s in $w_\bt$ for those moved from 
$w_\bk$ group by group. 
The first binomial factor $\binom{t_1}{k_1}$
of \eqref{ArrayBino} is the number of ways to choose~$k_1$ places for~$X$'s (coming from 
the first group in $w_\bk$) in the first group $X^{t_1}Y$ of $w_\bt$.
The second binomial factor $\binom{t_1+t_2-k_1}{k_2}$ of \eqref{ArrayBino} 
represents the number of ways to choose~$k_2$ 
places for~$X$'s (coming from the second group $YX^{k_2}Y$ in~$w_\bk$)
 in the first two groups $X^{t_1}YX^{t_2}Y$ of~$w_\bt$ 
where the already occupied $k_1$ places in the previous step
are prohibited from being chosen.
We continue the process in the same way. For each given 
$i\in\{2,\dots,d\}$, suppose that destinations of $X$'s 
in~$X^{t_1}Y\cdots X^{t_{i-1}}$ 
from $X^{k_1}Y\cdots YX^{k_{i-1}}Y$ have already been
chosen. Then, the $i$-th binomial
factor $\binom{t_1+\cdots t_i-k_1-\cdots-k_{i-1}}{k_i}$ 
of \eqref{ArrayBino}
represents the number of ways to choose $k_i$ places for 
$X$'s (coming from the $i$-th group $YX^{k_i}Y$ in~$w_\bk$)
in~$X^{t_1}Y\cdots YX^{t_i}Y$ (the first $i$ groups of~$w_\bt$):
there are $t_1+\cdots+t_i$ places for $X$ in~$X^{t_1}Y\cdots YX^{t_i}Y$,
but already $k_1+\cdots+k_{i-1}$ places are occupied 
by earlier choices.
Performing the process until $i=d$ 
verifies the desired identity
$\la w_\bt, \SS^{(\bk)} \ra=\binom{\bt}{\bk}$.
\end{Remark}


\begin{proof}[Proof of Theorem \ref{orthogonality}]
From the formula~$Y^{(k)}=\sum_{i=0}^k (-1)^i\binom{k}{i}X^{k-i}YX^i$~(\cite[(4)]{M37}), 
it is not difficult to see that the expansion of the Magnus polynomial in monomials
is given by 
\begin{equation} \label{Mag-Mono}
\MM^{(\bk)}=\sum_{\bt \in\N_0^{(\infty)}}
\magnus{\bk}{\bt} w_{\bt} 
\end{equation}
with
\begin{equation}
\label{MagnusBino}
\magnus{\bk}{\bt}:=(-1)^{\sum_{i=1}^d (d-i+1)(k_i-t_i)}
\binom{k_1}{k_1-t_1} 
\binom{k_2}{k_1+k_2-t_1-t_2} \cdots
\binom{k_d}{\sum_{i=1}^d(k_i-t_i)}
\end{equation}
for $\bt:=(t_1,\dots, t_d;t_\infty)$ and $\bk:=(k_1,\dots,k_d;k_\infty)$.
Since we have $\lp\SS^{(\bt)},\MM^{(\bk)}\rp=\sum_{\bu\in\N_0^{(\infty)}}
\lp\SS^{(\bt)},w_\bu\rp \lp\MM^{(\bk)},w_\bu\rp$, it suffices to show
\begin{equation} \label{target}
\sum_\bu \magnus{\bk}{\bu} \binom{\bu}{\bt}=\delta_\bk^\bt.
\end{equation}
Noting that a non-zero pairing $\lp\SS^{(\bt)},\MM^{(\bk)}\rp$ occurs only when
$|\bt|=|\bk|$ and $\dep(\bt)=\dep(\bk)$, we may assume without loss of generality 
that $\bu$ in the above summation also runs over those
with the fixed size $N:=|\bt|=|\bk|$ and depth $d:=\dep(\bt)=\dep(\bk)$.
Then, the summation $\sum_\bu$ with $\bu=(u_1,\dots,u_d;u_\infty)$ 
has $d$ independent parameters $u_1,\dots, u_d$ that determine 
$u_\infty=N-\sum_{i=1}^d u_i$. We may also regard each $u_i$ as running over $\Z$,
as the coefficients $\magnus{\bk}{\bu}$, $\binom{\bu}{\bt}$ vanish when 
their combinatorial meaning is lost.
Then, in the summation $\sum_{(u_1,\dots,u_d)\in\Z^d}$ in \eqref{target},
the partial factor of summation 
involved with the last parameter $u_d$ can be factored out in the form:
\begin{align*}
&
\sum_{u_d\in\Z}(-1)^{-u_d}
\binom{k_d}{u_d+\sum_{i=1}^{d-1}(u_i-k_i)}
\binom{u_d+\sum_{i=1}^{d-1}(u_i-t_i)}{t_d} \\
&=
(-1)^{\sum_{i=1}^{d-1}(u_i-k_i)-k_d}
\binom{\sum_{i=1}^{d-1}(k_i-t_i)}{t_d-k_d}.
\end{align*}
(Use \cite[(5.24)]{GKP}.)
Repeating this process inductively on $d$, we eventually find
$$
\lp\SS^{(\bt)},\MM^{(\bk)}\rp=
\binom{0}{t_1-k_1}
\binom{k_1-t_1}{t_2-k_2}
\binom{k_1+k_2-t_1-t_2}{t_3-k_3}
\cdots
\binom{\sum_{i=1}^{d-1}(k_i-t_i)}{t_d-k_d}
$$
which is equal to $\delta^\bk_\bt$ as desired.
\end{proof}


\begin{Corollary} \label{coro}
Each element $u\in R\la X,Y\ra$ can be written as 
$$
u=\sum_{\bk \in\N_0^{(\infty)}}
\lp\SS^{(\bk)}, u\rp \,\MM^{(\bk)}
=\sum_{\bk \in\N_0^{(\infty)}}\lp\MM^{(\bk)}, u\rp \,\SS^{(\bk)}.
$$
\end{Corollary}

Note that only a finite number of summands are nonzero
in either summation above. 
